\begin{afterword}
\summaryhead

The post quotes Robyn Dawes on Norman Maier's rule for groups: ``Do not propose solutions until the problem has been discussed as thoroughly as possible without suggesting any.'' In Maier's role-play, three assembly-line workers and a foreman discuss an expert's advice to stop rotating jobs. Groups given the rule more often found the better arrangement, in which the two abler workers rotate and the least able keeps the easiest job. Dawes adds that he uses the rule with his own groups, especially on tough problems.

The author agrees and says the problem grows with difficulty: many people claim to know how to build an artificial general intelligence, and a majority solve Friendly AI in fifteen seconds. The rule fits ``The Bottom Line'': a decision is only as good as the thinking done before it was made. Since once you can guess your answer you have probably decided, and evidence rarely dislodges an idea, the author suspects that a better method than falsification is to hold off on thinking of an answer at all.

\responsehead

The post rests on a rule with a real source, and it credits the source. What it adds, it states with more confidence than its evidence gives.

The evidence is one role-play, reported secondhand. Groups told to discuss the problem first more often found a better arrangement for three workers, a problem with a known best answer. Dawes says the rule is easy to show working where good solutions are objectively defined. For the tough problems of his own groups, he says he has no objective criterion, and the rule only ``appears'' to help. The post's response is ``This is so true it's not even funny.'' I could not reach Dawes's book or Maier's papers, so I cannot say how large Maier's effect was.

The post then moves the rule from groups to single minds. Maier's reason for the rule concerns groups: members become attached to the solutions they have proposed, and the discussion turns to those. The post's own method, to ``hold off on thinking of an answer,'' concerns one person's thinking. It is offered with ``I suspect,'' and nothing in the post tests it. It is also old. In \textsc{How We Think} (1910) John Dewey called suspended judgment ``the essence of critical thinking,'' and described it as inquiry into the nature of the problem ``before proceeding to attempts at its solution.''

The examples are the author's encounters with people who solve artificial intelligence at once, told as jokes. The footnote cites the author's own chapter, which reports the same kind of encounter and gives no count.

In short: a sound and properly credited rule for groups, affirmed with more certainty than its one role-play supports, then extended to individual thinking as a suspicion, untested and without mention of the older version of the idea.
\end{afterword}
